\subsection{Proof by Contraposition}
\label{subsections:techniques:universal:contraposition}



\paragraph{Description.}
A proof by contraposition proves a conditional statement by proving an %
equivalent statement instead. %
Precisely, it proves a universal conditional statement, %
``$\forall s \in S$, if $P(s)$, then $Q(s)$,'' %
by proving its contrapositive, %
``$\forall s \in S$, if $\neg Q(s)$, then $\neg P(s)$.'' %
Generally, a proof by contraposition mirrors the following format:
\begin{enumerate}{
    \item \textbf{Setup.} %
        Let $s \in S$ be arbitrary and assume that $\neg Q(s)$ is true.
    \item \textbf{Verification.} %
        Show that $\neg P(s)$ is true.
    \item \textbf{Conclusion.} %
        Conclude the proof, optionally by summarizing what has been found.
}
\end{enumerate}



\paragraph{Advice.}
\begin{itemize}{
    \item Use this proof for a conditional statement when assuming %
        $\neg Q(s)$ is easier than assuming $P(s)$, %
        such as when $Q(s)$ is a negation.
    \item State the contrapositive, the conditional with its hypothesis %
        and conclusion swapped and negated, at the start of the proof, %
        as assuming $\neg Q(s)$ can otherwise seem random.
}
\end{itemize}



\paragraph{Example.}
\begin{exercise}{
    Prove that $\forall z \in \Z$, if $z^2$ is even, then $z$ is even.
}
\end{exercise}
\begin{proof}{
    This statement can be proven by contraposition.

    \begin{annotatedpanel}[Setup.]{
        Consider the contrapositive: if $z$ is odd, then $z^2$ is odd. %
        Let $z \in \Z$ such that $z$ is odd. %
        Then $z = 2k + 1$ with $k \in \Z$.
    }{
        \textbf{Step 1.} Let $s \in S$ be arbitrary and assume that %
        $\neg Q(s)$ is true. %
        Here, $S = \Z$, $s = z$, $P(z)$ is the statement that $z^2$ is even, %
        $Q(z)$ is the statement that $z$ is even, %
        and so $\neg Q(z)$ is the statement that $z$ is odd.
    }
    \end{annotatedpanel}

    \begin{annotatedpanel}[Verification.]{
        Then %
        \begin{align*}
            z^2 &= (2k + 1)^2\\
                &= 4k^2 + 4k + 1\\
                &= 2(2k^2 + 2k) + 1.
        \end{align*}
        Let $m \in \Z$ such that $m = 2k^2 + 2k$. %
        Then $z^2 = 2(2k^2 + 2k) + 1 = 2m + 1$. %
        Thus, $z^2$ is odd.
    }{
        \textbf{Step 2.} Show that $\neg P(s)$ is true. %
        Here, $\neg P(z)$ is the statement that $z^2$ is odd, %
        which holds because $z^2 = 2m + 1$.
    }
    \end{annotatedpanel}

    \begin{annotatedpanel}[Conclusion.]{
        Thus, $\forall z \in \Z$, if $z$ is odd, then $z^2$ is odd, %
        so the contrapositive holds. %
        Therefore, the statement holds by contraposition. \qedhere
    }{
        \textbf{Step 3.} Conclude the proof. %
        Here, the summary states what has been found: %
        the contrapositive, if $\neg Q(z)$, then $\neg P(z)$, holds.
    }
    \end{annotatedpanel}
}
\end{proof}